% Generated draft: requires mathematical and source review.
\documentclass[11pt]{article}
\usepackage[margin=25mm]{geometry}
\usepackage{fontspec}
\setmainfont{Latin Modern Roman}
\usepackage{amsmath,amssymb,mathtools,unicode-math}
\setmathfont{Latin Modern Math}
\usepackage{enumitem,needspace,xurl,hyperref}
\setlength{\parindent}{0pt}
\setlength{\parskip}{5pt}
\setlength{\emergencystretch}{3em}
\newcommand{\N}{\mathbb{N}}
\newcommand{\Z}{\mathbb{Z}}
\newcommand{\Q}{\mathbb{Q}}
\newcommand{\R}{\mathbb{R}}
\newcommand{\C}{\mathbb{C}}
\newcommand{\F}{\mathbb{F}}
\newcommand{\A}{\mathbb{A}}
\newcommand{\PP}{\mathbb{P}}
\newcommand{\E}{\mathbb{E}}
\newcommand{\eps}{\varepsilon}
\newcommand{\dd}{\,\mathrm{d}}
\begin{document}
\section*{UnsolvedMath ID: 30004654}
\textbf{Source problem code:} catalogue entry
\subsection{Definitions and mathematical statement}
For a fixed knot model, let \(\mathcal K\) be its set of isotopy classes of classical knots in \(\R^3\), and let \(\Z[\mathcal K]\) be the free abelian group on this set. The source does not specify whether knots are long or closed, oriented or unoriented, or framed or unframed; these choices must be fixed consistently before specifying the embedding-calculus tower.

For \(r\geq 1\), let \(F_r\subseteq\Z[\mathcal K]\) be generated by the alternating resolution sums
\[
\sum_{\varepsilon\in\{+,-\}^{r}}(-1)^{\#\{i:\varepsilon_i=-\}}[K_{\varepsilon}],
\]
where a singular knot has exactly \(r\) transverse double points and \(K_{\varepsilon}\) denotes the knot obtained by resolving each double point with the indicated crossing sign. For an abelian group \(A\), an invariant \(v:\mathcal K\to A\) has finite type at most \(m\) if its linear extension to \(\Z[\mathcal K]\) vanishes on \(F_{m+1}\). Thus
\[
U_m=\Z[\mathcal K]/F_{m+1},\qquad u_m(K)=[K]\bmod F_{m+1},
\]
is a universal integral invariant of type at most \(m\): every such \(A\)-valued invariant factors uniquely through a group homomorphism \(U_m\to A\).

Let \(\mathscr K\) denote the chosen space of knots, and write
\[
\eta_n:\mathscr K\longrightarrow T_n\mathscr K
\]
for its map to the \(n\)th Goodwillie--Weiss approximation, once the knot model and tower convention are fixed. The supplied question asks whether, for each applicable positive integer \(n\), this approximation gives a universal integral finite-type invariant of order at most \(n-1\), without inverting any primes. In the universal abelian-group formulation above, this asks whether the invariant extracted from the stage has the universal property of \(u_{n-1}\), with coefficients over \(\Z\).

\textit{Unresolved formulation qualifications.} The source does not specify the invariant extracted from \(T_n\mathscr K\), its target and algebraic structure, the precise meaning of universality, or the admissible low values of \(n\). Consequently, the universal abelian-group formulation explains a standard meaning of integral universality, but its identification with the intended tower-stage assertion requires review. No identification of \(\pi_0(T_n\mathscr K)\) with \(U_{n-1}\) is asserted here.
\subsection{Short English statement}
The question asks whether the \(n\)th Goodwillie--Weiss approximation captures universally all finite-type knot invariants of order at most \(n-1\) over the integers, without inverting primes. A precise tower-stage statement still requires the knot model, indexing convention, target invariant, and meaning of universality to be specified.
\subsection{Sources}
\begin{itemize}
\item Local registry: supplied source record for permanent UnsolvedMath ID \(30004654\), source problem code \texttt{catalogue entry}, titled \emph{Integral Universality of Goodwillie-Weiss Knot Invariants}. The recorded proposed year is 2021; no proposer is named.
\item Supplied report citation: \emph{Homotopic and Geometric Galois Theory} (2022), Oberwolfach Reports, pp.~663--744, DOI \texttt{10.4171/owr/2021/12}. \url{https://doi.org/10.4171/owr/2021/12}. The background describes Oberwolfach Reports as documenting mathematical workshops and their open-problem discussions. The supplied literature list labels this reference \emph{Oberwolfach Report 12/2021}.
\item Supplied literature reference: Turchin--Willwacher, \emph{Embedding calculus knot invariants}. \url{https://arxiv.org/abs/1411.1832}.
\item Supplied literature reference: Goodwillie--Klein--Weiss, \emph{embedding calculus and grope cobordism}. \url{https://arxiv.org/abs/2010.05120}.
\item Literature assessment dated exactly \texttt{2026-08-22}: the supplied curation labels the problem \texttt{partially\_solved} and states that Goodwillie--Weiss stages are universal after localization at suitable primes, that finite-type calculations support the integral conjecture, and that full integral universality is not proved. These are supplied dated curation claims, not independently verified mathematical or literature conclusions. The source explicitly requests independent verification of time-sensitive status and cited claims.
\item The registry labels the statement \texttt{corrected\_verified} and records a manual semantic, typographic, and LaTeX review claiming that its clean statement is self-contained and preserves the target. This is supplied review metadata; it does not resolve the formulation qualifications identified here.
\end{itemize}
\end{document}
